\section{ماژول mamba\_shedder (Mamba-Shedder Compressor)}
\label{sec:module_mamba_shedder}

\subsection{معرفی و انگیزه}
این مکانیزم بر مبنای مبانی نظری پیشرفته در ادبیات توجه و فشرده‌سازی \cite{gu2023mamba} توسعه یافته است.

افزایش عمق شبکه‌های زبانی متوالی و ترنسفورمرها سبب تحمیل سربار محاسباتی سنگین در استنتاج بلادرنگ می‌گردد. در معماری‌های عمیق، بسیاری از لایه‌های متوالی تکراری دارای اطلاعات تفکیکی حاشیه‌ای بوده و یا کانال‌های ورودی آن‌ها مشارکت ناچیزی در تولید خروجی ایفا می‌کنند. کاهش وزن‌ها به تنهایی برای مهار حافظه و تأخیر کافی نیست؛ بلکه نیاز به هرس ساختاریافته در دو سطح وجود دارد: هرس کانال‌های کم‌اهمیت (\lr{Channel Pruning}) و حذف کامل لایه‌های اضافی (\lr{Layer Shedding}).

فشرده‌ساز \lr{Mamba-Shedder} راهکاری یکپارچه برای هرس ساختاریافته مبتنی بر حساسیت اطلاعات فیشر (\lr{Empirical Fisher Information}) و حذف گزینشی لایه‌ها ارائه می‌دهد. در این الگو، میزان وابستگی خروجی به هر کانال ورودی با تخمین ماتریس اطلاع فیشر تجربی ارزیابی شده و کانال‌های کم‌انرژی هرس می‌شوند. هم‌زمان، لایه‌های با حساسیت باقیمانده اندک از ساختار مدل جدا شده و عمق کلی شبکه تقلیل می‌یابد که به شتاب‌بخشی محسوس در استنتاج می‌انجامد.

\subsection{مبانی نظری و فرمول‌بندی ریاضی}
فرمول‌بندی ریاضی الگوریتم \lr{Mamba-Shedder} مبتنی بر دو مکانیزم تحلیلی است:

\begin{enumerate}
    \item \textbf{تخمین حساسیت بر پایه اطلاعات فیشر تجربی (\lr{Empirical Fisher Sensitivity}):} برای هر کانال ورودی $j \in \{1, \dots, d_{\text{in}}\}$ در لایه خطی با ماتریس وزن $\mathbf{W} \in \mathbb{R}^{d_{\text{out}} \times d_{\text{in}}}$ و تانسور فعال‌سازی ورودی کالیبراسیون $\mathbf{X} \in \mathbb{R}^{N \times d_{\text{in}}}$، شاخص حساسیت اطلاعات فیشر تجربی به شکل زیر تعریف می‌شود:
    \begin{equation}
    \mathcal{I}(X_j) = \sqrt{\frac{1}{N} \sum_{n=1}^N X_{n, j}^2 + \epsilon}
    \end{equation}
    در صورت عدم دسترسی به داده‌های کالیبراسیون، هنجار ستونی ماتریس وزن $\|\mathbf{W}_{:, j}\|_2$ به عنوان شاخص حساسیت جایگزین مورد استفاده قرار می‌گیرد.
    
    سپس بر اساس درصد پراکندگی هدف $s \in [0, 1)$، ماسک دودویی کانال‌ها با حفظ $k = \lfloor d_{\text{in}}(1 - s) \rfloor$ کانال برتر تولید شده و ماتریس وزن هرس می‌شود:
    \begin{equation}
    \mathbf{W}_{\text{pruned}} = \mathbf{W} \odot \mathbf{M}_{\text{channel}}^\top
    \end{equation}
    همچنین ماژول توانایی هرس ساختاریافته ۲:۴ (\lr{NVIDIA 2:4 Structured Sparsity}) را داراست که در آن از هر ۴ عنصر مجاور در امتداد بعد ورودی، دقیقا ۲ عنصر با کمترین میزان شایستگی صفر می‌شوند.

    \item \textbf{شاخص حساسیت تفاضلی لایه‌ها و حذف لایه (\lr{Layer Shedding}):} حساسیت لایه $l$-ام بر اساس تغییرات تانسور بازنمایی در گذر از لایه محاسبه می‌گردد:
    \begin{equation}
    \Delta_l = \frac{\|\mathbf{X}_{l+1} - \mathbf{X}_l\|_F}{\|\mathbf{X}_l\|_F + \epsilon}
    \end{equation}
    در فرآیند حذف لایه، تعداد $N_{\text{shed}}$ لایه از میان بلوک‌های میانی شبکه بدون دستکاری لایه ابتدایی و انتهایی انتخاب و از کانتینر \lr{nn.ModuleList} حذف می‌شوند تا اتصال مستقیم میان لایه‌های باقیمانده برقرار گردد.
\end{enumerate}

\subsection{معماری و پیاده‌سازی}
کلاس \lr{MambaShedderCompressor} در ماژول \lr{Codes/framework/compression\_engine/mamba\_shedder.py} پیاده‌سازی شده و با شناسه‌های \lr{mamba\_shedder}، \lr{mamba\_shedder\_compressor} و \lr{mamba\_shedder\_factory} در رجیستری مرکزی قابل دسترس است. ویژگی‌های طراحی عبارتند از:
\begin{itemize}
    \item \textbf{مکانیزم حذف ایمن لایه‌ها با اعتبارسنجی خودکار (\lr{Layer Shedding with Revert}):} پس از حذف لایه‌های منتخب از مدل، یک گذر پیشرو آزمایشی اجرا می‌شود. در صورت بروز هرگونه ناهماهنگی در ابعاد پنهان یا عدم تطابق تانسورهای خروجی، ساختار اولیه لایه‌ها بازگردانی شده و از شکست مدل جلوگیری می‌شود.
    \item \textbf{پشتیبانی از ساختارهای گوناگون انکودر و دیکودر:} ماژول کانتینرهای مختلف لایه‌ها در مدل‌های استاندارد نظیر \lr{transformer.layer}، \lr{transformer.h}، \lr{model.layers}، \lr{encoder.layer} و \lr{backbone.layers} را به صورت خودکار شناسایی و ردیابی می‌کند.
    \item \textbf{پرهیز از هرس لایه‌های حساس:} ماژول به صورت قطعی لایه‌های طبقه‌بندی (\lr{Classifier})، سر نگاشت زبانی (\lr{lm\_head})، ماتریس‌های تعبیه (\lr{Embeddings}) و لایه‌های نرمال‌سازی را از هرس معاف می‌دارد.
    \item \textbf{پشتیبانی منعطف از شتاب‌دهنده‌های ۲:۴:} فعال‌سازی پرچم \lr{structured\_24} امکان بهره‌برداری مستقیم از موتورهای محاسباتی متراکم تنسور کُر (\lr{Tensor Cores}) انویدیا را مهیا می‌سازد.
\end{itemize}

\subsection{نتایج تجربی}
ارزیابی ماژول \lr{mamba\_shedder} بر روی مدل‌های معیار طبق ارقام مندرج در فایل تاییدشده \lr{thesis\_final/master\_verified\_results.json} در جدول \ref{tab:benchmark_mamba_shedder} منعکس شده است:

\begin{table}[H]
\centering
\small
\caption{نتایج تجربی و اعتبارسنجی بنچ‌مارک فشرده‌ساز \lr{mamba\_shedder} روی مدل‌های پایه}
\label{tab:benchmark_mamba_shedder}
\begin{tabular}{lccccc}
\toprule
\textbf{مدل پایه} \lr{(Base Model)} & \textbf{دقت} \lr{(Accuracy)} & \textbf{شاخص} \lr{(F1-Score)} & \textbf{حافظه} \lr{(Peak VRAM)} & \textbf{توان} \lr{(Throughput)} & \textbf{وضعیت} \\
\midrule
\lr{DistilBERT} & \lr{87.50\%} & \lr{0.8868} & \lr{337.59 MB} & \lr{180,820.54 tok/s} & \lr{PASS} \\
\lr{RoBERTa} & \lr{85.44\%} & \lr{0.8542} & \lr{361.54 MB} & \lr{120,068.28 tok/s} & \lr{PASS} \\
\lr{GPT-2} & \lr{56.88\%} & \lr{0.6835} & \lr{634.91 MB} & \lr{67,770.43 tok/s} & \lr{PASS} \\
\lr{Qwen} & \lr{51.49\%} & \lr{---} & \lr{---} & \lr{---} & \lr{PASS} \\
\bottomrule
\end{tabular}
\end{table}

نتایج فوق بازدهی برجسته الگوریتم را در انکودرهای دومنظوره به نمایش می‌گذارد؛ در مدل \lr{DistilBERT} دقت مطلوب \lr{87.50\%} همراه با شاخص F1 برابر \lr{0.8868} و توان پردازشی چشمگیر \lr{180,820.54 tok/s} ثبت شده است. در مدل \lr{RoBERTa} دقت \lr{85.44\%} با توان \lr{120,068.28 tok/s} حاصل گردیده است. هم‌چنین مدل بزرگ \lr{Qwen} به دقت \lr{51.49\%} و \lr{GPT-2} به دقت \lr{56.88\%} دست یافته و همگی معیار قبولی (\lr{PASS}) را احراز کرده‌اند.

\subsection{محدودیت‌ها و تحلیل شکست}
با وجود انعطاف‌پذیری روش، محدودیت‌های فنی زیر قابل تامل است:
\begin{itemize}
    \item \textbf{تخریب مسیر انتشار گرادیان با حذف لایه‌ها در مدل‌های با عمق کم:} در مدل‌های فشرده نظیر \lr{DistilBERT} که دارای تنها ۶ لایه هستند، حذف حتی یک لایه منجر به از دست رفتن بیش از ۱۶ درصد ظرفیت پردازشی مدل شده و می‌تواند رانش بازنمایی جبران‌ناپذیری ایجاد کند.
    \item \textbf{انباشت خطای صفر کردن کانال‌ها در شبکه‌های دیکودر:} در مدل‌های علّی نظیر \lr{GPT-2}، صفر شدن کانال‌های ورودی لایه توجه بدون تغییر متناظر در لایه خروجی می‌تواند تعادل وزن‌های پیش‌بینی توکن‌های بعدی را بر هم زند و افت نسبی دقت را به همراه داشته باشد.
    \item \textbf{نیاز به کالیبراسیون داده‌محور جهت انتخاب بهینه لایه‌ها:} انتخاب لایه‌ها برای حذف صرفاً بر اساس الگوی تصادفی میانی نیازمند ارزیابی دقیق حساسیت باقیمانده بر روی دادگان ارزیابی است تا از حذف تصادفی لایه‌های حاوی اطلاعات ساختاری زبانی پیشگیری شود.
\end{itemize}
